\centerline{\bf \Huge Пересчитаем ребра I}

\begin{flushright}
        \scriptsize{}
\end{flushright}

\textbf{Определения.}
Будем называть \textit{графом} множество точек (\textit{вершин}), некоторые из которых соединены между собой линиями (\textit{рёбрами}). Любые две вершины могут быть соединены не более, чем одним ребром.

Граф называется связным, если от любой его вершины можно по рёбрам добраться до любой другой.

Если граф \textit{несвязен}, то он разбивается на части, которые называются \textit{компонентами связности}. Внутри каждой компоненты связности от каждой вершины существует путь по рёбрам до любой другой вершины. Между вершинами принадлежащими разным компонентам связности такого пути нет.

\textit{Степенью вершины} называется количество выходящих из этой вершины рёбер. Сумма степеней вершин любой компоненты связности равна удвоенному количеству рёбер.

\problem{1}
В стране 1568 городов, из каждого выходит по 4 дороги. Сколько всего дорог в стране?

\problem{2}
Можно ли изобразить на плоскости 9 отрезков так, чтобы каждый пересекал ровно 3 других? (считая, что в одной точке могут пересекаться не более двух отрезков)

\problem{3}
В некоторой стране 15 городов, каждый из которых соединен дорогами не менее чем с 7 другими. Докажите, что из любого города можно добраться до любого другого (возможно, проезжая через другие города).

\problem{4}
Каждый мальчик знаком с 7 мальчиками и 10 девочками, а каждая девочка — с 8 девочками и 9 мальчиками. Кого больше — мальчиков или девочек?

\problem{5}
У Элианы 5 друзей среди одноклассников. У остальных её одноклассников 4,6 или 8 друзей. И только у Эрдема всего один друг. Докажите, что Элиана может отправить Эрдему записку, если каждый будет передавать записку одному из своих друзей.

\problem{6}
Юноши и девушки пожимали друг другу руки (юноши — только девушкам, девушки — только юношам). Каждый юноша пожал 13 рук, а каждая девушка — 10 рук. При этом какие-то пары могли жать руки несколько раз. Докажите, что было не менее 130 рукопожатий.